\documentclass[11pt]{article}
% cmap + T1 + a Type 1 text font so fi/fl ligatures survive copy-paste and
% full-text indexing. Without these, "five" extracts as " ve".
\usepackage{cmap}
\usepackage[T1]{fontenc}
\usepackage{lmodern}
\usepackage[utf8]{inputenc}
\usepackage[margin=1in]{geometry}
\usepackage{graphicx}
\usepackage{booktabs}
\usepackage{amsmath}
\usepackage[hidelinks]{hyperref}
\usepackage{xurl}
\usepackage{xcolor}
\usepackage{enumitem}
\usepackage{caption}
\usepackage{subcaption}
\usepackage{float}
\usepackage{listings}
\usepackage{tikz}
\usetikzlibrary{shapes.geometric, positioning, arrows.meta, calc}
\usepackage{pgfplots}
\pgfplotsset{compat=1.18}
\usepackage{natbib}
\usepackage{needspace}
\usepackage{titlesec}
\newcommand{\sectionbreak}{\needspace{6\baselineskip}}
\bibliographystyle{unsrtnat}

\setlist{nosep,leftmargin=*}
\setlength{\parskip}{0.4em plus 0.1em minus 0.1em}

\lstset{
  basicstyle=\small\ttfamily,
  breaklines=true,
  frame=single,
  backgroundcolor=\color{gray!5},
  rulecolor=\color{gray!40},
  columns=fullflexible,
  keepspaces=true,
}

\definecolor{realgreen}{HTML}{2E7D32}
\definecolor{synthblue}{HTML}{1565C0}

\title{\textbf{IVF-Bench: A Rubric-Based Standard for Evaluating\\Vision-Language Models on IVF Clinical Reasoning}}
\author{
  Andrew G. A. Correa\\Monostate\\\texttt{andrew@monostate.ai}
  \and
  Brittany Yoon\\The Fertility Plan\\\texttt{b@thefertilityplan.com}
}
\date{August 2026}

\begin{document}
\maketitle

\begin{abstract}
ABSTRACT_PLACEHOLDER
\end{abstract}

% ============================================================
\section{Introduction}

Few decisions in IVF carry more weight than the choice of which embryo to
transfer, and in most clinics that choice still rests on visual assessment: an
embryologist examines the blastocyst on day 5 and assigns a Gardner grade
capturing the degree of expansion, the organization of the inner cell mass, and
the cohesion of the trophectoderm \citep{gardner1999}. The grade is clinically
useful, but it is also subjective, and when embryologists at different centers
are shown the same day-5 embryos and asked to select one for transfer, their
agreement falls well short of what the confidence of the decision would
suggest \citep{storr2017}.

That variability has motivated a decade of machine learning work, nearly all of
it framed as a scoring problem in which an image or a time-lapse sequence is
mapped to a scalar and embryos are ranked by that scalar. The largest randomized
evaluation of this paradigm, a double-blind trial across 14 clinics and 1,066
patients, compared a deep learning ranker against trained embryologists and did
not demonstrate noninferiority \citep{illingworth2024}, which suggests that
further refinement of image-only ranking is unlikely to be the most productive
direction.

Figure~\ref{fig:pipeline} contrasts that regime with the one this paper
proposes.

FIG_PIPELINE_PLACEHOLDER

We argue the limitation lies in the framing rather than the architecture. The
image captures only part of what determines whether a transfer succeeds; the rest
is patient-level, comprising ovarian reserve, endometrial receptivity, body mass
index and lifestyle, previous cycles, and the andrology profile. Most of that
resides with the treating physician and rarely reaches the embryologist at the
microscope, still less the imaging model, so a system operating on pixels alone
is not failing at prediction so much as being asked to predict without the
covariates that govern the outcome.

Vision-language models can in principle consume both modalities, but the claim
has never been testable, because embryo images are sensitive clinical records,
outcome-linked patient context more sensitive still, and the review required to
assemble the two is stringent enough that essentially every substantial dataset
in this field stays private. That, rather than model capacity, is the binding
constraint, and it motivates an instrument for measuring whether a model can use
patient context at all.

\textbf{IVF-Bench} comprises 753 blastocyst cases, each imaged on the day the
embryo reached the blastocyst stage and was transferred, 641 on day 5 and 112 on
day 4, and each pairing a real embryo image and its silver-standard Gardner
annotation with the cycle data and recorded outcomes of the same patient, together with the patient-history fields that no
public dataset supplies and that we therefore generate from published population
distributions, labeling each as generated and defending the choice in
Section~\ref{sec:synthetic}. A model is asked to describe the embryo, integrate
the patient factors, estimate an implantation probability, and recommend next
steps, and its answer is scored on five clinical rubrics while the stated
probability is compared against the recorded outcome. The generated fields are
sampled without reference to those outcomes and add no prognostic information
beyond the measured covariates, as Section~\ref{sec:synthetic} quantifies, so
every outcome-prediction result reported here is bounded by the data before any
model is involved.

We evaluate seven frontier and open-weight systems and then post-train a 9B open
model on preference pairs derived from the benchmark, obtaining the following results:

\begin{enumerate}
\item \textbf{Five hundred preference pairs raise a 9B open model 24.6\% over its
      own base on the held-out split}, and a judge from a different family
      independently puts the same gain at 15.5\%. This is the paper's central
      training result, a within-model comparison no leaderboard convention can
      distort, and it says the reasoning half of embryo assessment is reachable
      far below frontier scale. The gain is also enough to place the model second
      of eight ahead of Claude Opus 4.6 at roughly one seventh of the marginal
      inference cost, but that ranking is GPT-5.4's and reverses under Sonnet 4.6,
      so we report it as a consequence rather than as the claim.
\item The improvement spreads across all five rubrics rather than concentrating
      in one, which separates a procedure that raises the assessment from one
      that optimizes a reported quantity.
\item Outcome discrimination stays weak for every system, and we quantify the
      deficit against a baseline that ignores its inputs entirely, placing the
      ceiling in the data rather than in model capacity.
\item Morphology grounding is the weakest of the five rubrics, roughly 1.4 points
      below clinical integration, and putting the judge in front of the embryo is
      what revealed it: the same responses scored half a point higher under a
      judge that could not check them.
\end{enumerate}

Building the instrument surfaced four properties of rubric-based judging that
affect any benchmark of this shape, one of which withdraws a quantitative claim
the earlier release carried. They are incidental to the argument above and are
collected in Section~\ref{sec:judgelessons}.

The practical implication is that progress in this area depends less on model
development than on the construction of a public dataset linking embryo images to
measured patient context and recorded outcomes, and IVF-Bench provides the
instrument with which the value of such a dataset can be quantified.

% ============================================================
\section{Related work}

\textbf{Embryo selection models.} Most published systems predict a label from
pixels, whether by grading blastocysts automatically, ranking embryos within a
cohort, or estimating implantation from time-lapse sequences, and the commercial
products deployed in clinics follow the same pattern, positioned as adjunctive
scoring tools rather than autonomous decision makers; the strongest randomized
evidence for any of them is the trial discussed above \citep{illingworth2024}.
None of this work evaluates open-ended clinical reasoning, for the simple reason
that none of it produces language.

\textbf{Language about embryos.} The group that released the dataset we build on
has recently moved in this direction, publishing an expert-annotated set of
embryo images paired with natural-language morphological descriptions
\citep{neu2026dataset} together with a fine-tuned PaliGemma-2 model that generates
such descriptions \citep{neu2026invitrovision}. That work and ours are
complementary but differ in what they ask of a model: they evaluate whether a
model can describe an embryo, whereas we evaluate whether it can reason about the
patient carrying that embryo, which additionally requires the cycle data, the
history, a calibrated probability, and a recommendation. Their descriptions are
the morphology ground truth our evaluation most lacks, as we note in
Section~\ref{sec:limitations}.

\textbf{Rubric-graded evaluation of medical language models.} The practice of
scoring open-ended clinical answers against written criteria, with a language
model performing the scoring, was established at scale by HealthBench, for which
262 physicians authored 48,562 rubric criteria across 5,000 conversations
\citep{arora2025healthbench}. Our five rubrics are considerably narrower, written
for a single decision rather than for medicine in general, but the methodological
inheritance is direct, and we adopt the same principle of grading against
criteria rather than against a reference answer.

\textbf{Preference optimization.} We post-train with ORPO
\citep{hong2024orpo}, which folds the preference objective into the supervised
loss and removes the need for a separate reference model, in the family of
methods opened up by DPO \citep{rafailov2023dpo}.

% ============================================================
\section{The benchmark}

\subsection{Where the data comes from}

Open embryo data is rare, because embryo images are sensitive clinical records
and most published systems are trained on private single-clinic collections. We
build on the Kromp blastocyst dataset \citep{kromp2023}, which is public under CC
BY 4.0 and links blastocyst images to Gardner grades, selected cycle variables,
and pregnancy outcomes. Culture ran to the blastocyst stage and transfer followed
on day 4 where blastocysts were already available and on day 5 otherwise, so each
image is taken on its own case's transfer day rather than on a fixed day for
every case; the dataset records that day and we pass it to the model. The clinical annotation file contains 754 records spanning 732 patients, one of
which references an image absent from the distributed archive, leaving 753 usable
cases from 731 patients.

Preparing it took more correction than we expected, and we document that because
anyone reusing the dataset will meet the same issues. The annotation file is
semicolon-separated, Latin-1 encoded, and had at some point been opened in a
German-locale spreadsheet, which silently turned 501 AMH values and 144
endometrial thickness values into dates, so that 1.26 appears as
\texttt{J\"an.26}; we map the month abbreviations back. Oocyte counts sometimes
appear additively, as \texttt{7+1}, which we sum, and missing values stay
missing, shown to models as ``not available'' rather than imputed. One record
carries a live birth without the clinical pregnancy that must precede it, and we
leave it unaltered, since silently repairing a source dataset is less defensible
than reporting the inconsistency.

\subsection{What a case contains}

Each case combines measured fields drawn from Kromp with generated patient
context: the measured fields are the embryo image, the Gardner grades, patient
age, AMH, endometrial thickness, oocytes retrieved, mature oocytes, transfer day,
and the three pregnancy outcomes, while the generated fields are body mass index,
basal FSH, diagnosis, stimulation protocol, previous cycles and their outcomes,
partner age, semen parameters, and a lifestyle summary.

\subsection{Why the patient context is generated, and what that costs}
\label{sec:synthetic}

An embryo image accompanied by an age does not constitute a case: a clinician
reasoning about a transfer draws on the patient's history, and a benchmark that
omits that history cannot determine whether a model uses one. No public dataset,
however, carries body mass index, lifestyle, stimulation protocol, or cycle
history alongside embryo images and outcomes. Such data exist in clinics, largely
in physicians' notes and frequently only in their recollection, and they do not
circulate: they seldom reach the embryologist at the microscope and have never
reached a public research corpus, because assembling images together with
outcome-linked patient records is precisely the combination that regulatory and
privacy review is designed to constrain.

We therefore generate these fields, sampling each from a published population
distribution, with every parameter and its source given in
Appendix~\ref{app:distributions}. Generation is deterministic per case, seeded
from a hash of the case identifier, and each case record enumerates its generated
fields explicitly, so that no result can depend on generated values being
mistaken for measured ones. Identifiers are assigned positionally over the sorted
corpus, so regeneration is exact for a fixed corpus, while changing its size or
the split proportions reassigns identifiers and therefore resamples the profiles
that depend on them.

This procedure yields a case-level reasoning task that can be constructed and
reproduced entirely from public sources, at two costs that bear directly on the
interpretation of our results. First, the generated fields are sampled without reference to the recorded
outcomes. Two of them are nonetheless conditioned on the patient's real age,
namely basal FSH and partner age, which correlate with it at 0.32 and 0.73
respectively, and each therefore functions as a weak age proxy: taken alone, and
read in the direction that makes them predictive at all, with higher values
indicating a poorer prognosis, they reach AUROC 0.527 and 0.532 against clinical
pregnancy. Because age is itself
supplied as a measured field, these fields add no information a model did not
already have, but the stronger claim that the generated context is entirely
uninformative is not one we can make. The outcome metrics in
Section~\ref{sec:results} should be read as a floor that quantifies what is
missing rather than as a verdict on what vision-language models can achieve given
measured context. Second, apart from those two dependencies and the conditioning
of previous outcomes on previous cycle count, the fields are sampled
independently of one another, so genuine clinical correlations, such as that
between polycystic ovary syndrome and body mass index, are absent, and a model
that has internalized those correlations receives no credit for them here.
Neither limitation can be remedied by a more sophisticated generator; both are
remedied by a real dataset, which is precisely the argument of this paper.

\subsection{Ethics and data use}
\label{sec:ethics}

This work is a secondary analysis of a publicly released, de-identified dataset
and required no additional review. The underlying embryo images and clinical
records were collected under approval from the Ethics Committee of the Faculty of
Medicine at Johannes Kepler University Linz (Nr.\ 1238/2021) with informed
consent obtained from participants, and were released by the original
investigators under CC BY 4.0 \citep{kromp2023}; we redistribute no images and
our release regenerates every case from that source. The patient-history fields
we add are sampled from published population distributions and correspond to no
real individual.

The model released with this benchmark is a research artifact and is not a
medical device. It has no regulatory clearance, has not been evaluated in a
clinical trial, and must not be used to inform the care of any patient. We regard
the risk of misuse as material rather than hypothetical, because the model
produces fluent, confident, guideline-shaped counseling text whose surface
quality substantially exceeds its demonstrated prognostic validity, which is the
combination most likely to induce unwarranted confidence in a reader. The
strongest randomized evidence in this area found that a deep learning system did
not outperform trained embryologists \citep{illingworth2024}, and nothing in this
paper meets a lower bar than that.

\subsection{Splits}

The 753 cases are partitioned deterministically into 550 test, 100 validation,
and 103 held out, with the test split carrying the main leaderboard. The
validation split was reserved for judge and threshold calibration and, in the end,
was not used: no system was run against it, and we release it unused so that
future work has a split untouched by anything reported here. The held-out split
remained untouched through all training, hyperparameter search, and judge
calibration, and exists to answer a single question: whether the test-split
ordering survives on cases that nothing was optimized against.

One property of these splits should be stated plainly, since it is a defect: they
were drawn per image rather than per patient, and because twenty-two patients
contributed two embryos each, for six of those patients the two embryos fell into
different splits, so that four of the 103 held-out cases share a patient with a
test-split case. We quantify the effect in Section~\ref{sec:robustness}, where it
changes no ranking, but a subsequent version should partition on patient.

\subsection{The prompt}

Every model receives the embryo image together with a structured case
description containing the Gardner grade with plain-language descriptors, the
laboratory fields, and the patient profile, and is asked for five components: a
morphology description, an integration of the patient factors, the risks and
favorable features, a numeric implantation probability, and recommendations for
the present transfer and for future cycles. The full template appears in
Appendix~\ref{app:prompt}.

\subsection{The five rubrics}

A judge model reads the case and the response and returns a 1 to 5 score with a
short justification for each of five rubrics.

\begin{itemize}
\item \textbf{Morphological accuracy.} Does the description of the embryo match the
      image and the Gardner annotation?
\item \textbf{Clinical integration.} Does the response connect the patient factors
      to this embryo and this prognosis, rather than listing them?
\item \textbf{Reasoning coherence.} Is the argument internally consistent, with
      uncertainty where uncertainty belongs?
\item \textbf{Guideline alignment.} Is the advice consistent with evidence-based
      reproductive medicine?
\item \textbf{Actionability.} Could a clinician act on this, or is it
      generic?
\end{itemize}

Full rubric text and the 1 to 5 anchors for each appear in
Appendix~\ref{app:rubrics}.

% ============================================================
\section{Evaluation setup}
\label{sec:setup}

RESULTS_SETUP_PLACEHOLDER

% ============================================================
\section{Results}
\label{sec:results}

RESULTS_PLACEHOLDER

% ============================================================
\section{Post-training a small model}
\label{sec:orpo}

ORPO_PLACEHOLDER

% ============================================================
\section{What the benchmark itself gets wrong}
\label{sec:robustness}

ROBUSTNESS_PLACEHOLDER

% ============================================================
\section{Four things we learned about grading with a language model}
\label{sec:judgelessons}

JUDGELESSONS_PLACEHOLDER

% ============================================================
\section{Limitations}
\label{sec:limitations}

LIMITATIONS_PLACEHOLDER

% ============================================================
\section{What we release}
\label{sec:release}

RELEASE_PLACEHOLDER

% no \clearpage: the references and appendix flow on rather than each claiming a
% fresh page, which was costing two nearly-empty pages
\bibliography{references}

\appendix
\section{Case prompt template}
\label{app:prompt}
PROMPT_APPENDIX_PLACEHOLDER

\section{Rubric definitions}
\label{app:rubrics}
RUBRIC_APPENDIX_PLACEHOLDER

\section{Generated patient-context distributions}
\label{app:distributions}
DISTRIBUTION_APPENDIX_PLACEHOLDER

\section{Three evaluation defects}
\label{app:defects}
DEFECTS_APPENDIX_PLACEHOLDER

\section{Code and data availability}
\label{app:artifacts}
ARTIFACT_APPENDIX_PLACEHOLDER

\end{document}
